% !TEX root = ../my-thesis.tex
%
\chapter{Introduction}
\label{sec:intro}

\section{Motivation}
\label{sec:intro:motivation}

Human argumentation is hard to capture by machines because it does not follow
the rules of classical logic. Instead of deriving truth from fixed premises, a
debate consists of arguments and counter-arguments, and evidence enters weighted
or by comparison with other evidence rather than as a pure fact that is true or
false. This resembles a verdict that rests on aggravating and mitigating
circumstances rather than on a single decisive proof. The problem is not narrow
or merely technical. It connects to deeper questions such as algorithmic
fairness, the reasoning behaviour of neural models, and the hallucinations of
language models.

One established way to formalise arguments is the logic of
argumentation~\cite{krause1995logic}, which provides a language for representing
the structure of reasons that support or undercut a line of reasoning. While
this approach is useful and valid, several open questions remain. In particular,
it is not clear where arguments and counter-arguments come from in practical
applications and how they are in fact related. The examples in the literature
often appear constructed, and providing the arguments frequently seems harder
than the reasoning over them.

Large language models (LLMs) sit at the opposite end. They can produce reasons
for or against a cause in fluent natural language, but they lack the formal
structure that would allow any systematic analysis of how those reasons relate.
A standalone LLM can be asked to summarise a debate, yet its answer is another
ungrounded generation without a guarantee of internal consistency. The proposal
of this thesis is to use the LLM for what it is good at, namely producing
arguments and identifying their attacks, and to hand the resulting structure to
a formal layer that is good at reasoning about consistency. The formal layer
reconstructs deeply layered argumentations and makes explicit which sets of
arguments can be jointly accepted, a property that the raw LLM output leaves
implicit.

\section{Problem Statement and Research Questions}
\label{sec:intro:rqs}

This thesis builds and evaluates an end-to-end pipeline that turns
natural-language debate text into a formally evaluated Abstract Argumentation
Framework (AAF) in the sense of Dung~\cite{dung1995}. An LLM proposes
categorized attack candidates between arguments, a confidence filter routes them
to acceptance, rejection, or verification, and a Prolog solver computes grounded,
preferred, and stable extensions over the accepted attacks.

The work is guided by four research questions.

\begin{description}
  \item[RQ1 (Feasibility).] Can an end-to-end pipeline be built that turns
    natural-language arguments into formal Dung extensions through LLM-based
    attack extraction and a symbolic solver, and what do the resulting
    frameworks look like across a range of debate topics?
  \item[RQ2 (Typed vs.\ binary extraction).] Does \emph{typed} attack prompting,
    which forces the model to name the attack type and justify it, produce
    different or more reliable attack relations than plain \emph{binary}
    prompting, and how does this affect the resulting frameworks?
  \item[RQ3 (Extraction quality).] How well do the extracted attacks match a
    manual gold annotation in terms of precision, recall, and F1, and which
    error types dominate?
  \item[RQ4 (Argumentative structure).] What argumentative structures do the
    formal semantics reveal across the eight UKP topics, in particular regarding
    cycles and multiple extensions, and when does the formal layer add
    information beyond the grounded extension?
\end{description}

\section{Contributions}
\label{sec:intro:contributions}

The thesis makes the following contributions.

\begin{itemize}
  \item A complete and reproducible pipeline that connects LLM-based argument
    mining with Dung's abstract argumentation, from atomization through
    categorized attack extraction and confidence filtering to a Prolog solver.
  \item An empirical analysis of the resulting frameworks over eight topics of
    the UKP Sentential Argument Mining Corpus~\cite{stab2018}, reporting
    framework size, cycles, and the number of grounded, preferred, and stable
    extensions (RQ1, RQ4).
  \item An ablation that compares typed against binary attack prompting on an
    identical argument set, isolating the effect of the categorization that the
    underlying paper advertises but does not measure (RQ2).
  \item A small manually annotated gold standard that yields precision, recall,
    and F1 for the attack detection on one topic (RQ3).
\end{itemize}

\noindent
A condensed version of the pipeline and the cross-topic evaluation has been
submitted to IJCLR. This thesis extends that work with explicit research
questions, an expanded background, and the two additional experiments for RQ2
and RQ3.

\section{Thesis Structure}
\label{sec:intro:structure}

\textbf{Chapter \ref{sec:foundations}} introduces the formal background, namely
Dung's abstract argumentation frameworks and their semantics, the attack
typology of Pollock and ASPIC+, and the basics of argument mining and language
models.

\textbf{Chapter \ref{sec:related}} reviews related work on LLM-based
argumentation pipelines and positions the thesis against it.

\textbf{Chapter \ref{sec:method}} describes the pipeline in detail, step by step,
together with the design decisions behind it.

\textbf{Chapter \ref{sec:evaluation}} presents the experimental setup and
results, including the cross-topic evaluation, the typed-versus-binary ablation,
the gold-standard evaluation, and a detailed case study.

\textbf{Chapter \ref{sec:discussion}} discusses what the formal layer adds,
the threats to validity, and the main design trade-offs.

\textbf{Chapter \ref{sec:conclusion}} summarises the answers to the research
questions and outlines future work.
